\begin{proofbox}
\paragraph{最低门槛的证明：差商与左右大小}
先证必要条件。设 $x_0$ 为局部极大值点且可导。充分小的 $h$ 满足
\[
\frac{f(x_0+h)-f(x_0)}h\le0\quad(h>0),\qquad
\frac{f(x_0+h)-f(x_0)}h\ge0\quad(h<0).
\]
左右极限都等于 $f'(x_0)$，所以它同时 $\le0$ 与 $\ge0$，即 $f'(x_0)=0$。
极小值将不等号反向，仍得同一结论。

再证变号判定。若左正右负，由单调性定理与 $x_0$ 处连续性，$f$ 在包含 $x_0$ 的左半段严格递增，在右半段严格递减。
因此对两侧足够近的 $x\ne x_0$，都有 $f(x)<f(x_0)$，即严格极大值。左负右正同理得到严格极小值。
若两侧均为正，则附近左侧值小于 $f(x_0)$，右侧值大于 $f(x_0)$，故既非极大也非极小；两侧均为负同理。

这说明求出 $f'=0$ 只是找到候选点，不能省略附近大小的检查。
\end{proofbox}
